\newcommand{\sporkGrammarStuff}{
    \footnotesize
    \centering
    \arraycolsep=1.8pt\def\arraystretch{1.0}
    \setlength{\tabcolsep}{2pt}
}
\newcommand{\sporkSemanticsStuff}{
    \centering
    \small
}

\newcommand{\textgpu}[1]{\textsf{#1}}
\definecolor{barriercolor}{HTML}{BB00BB}

%% "Highlight changes" form of small step semantic rules.
%% In case of haters, redefine macros so they generate normal verbose rules.
%% i.e. redefine so \sporkRule generates the \langle ... \rangle \to \langle ... \rangle
%% rule with the LHS/RHS substitution of each \sporkSub.
%% Intentionally subdued colors (they are a pure scanning convenience; \blacktriangleright clearly separates).
%% Reserve more bold colors for more important info.
%% However, contrast with white background is important
%% as the bounds of the surrounding box affect the meaning.
\definecolor{sporkSubLhsColor}{HTML}{000000}
\definecolor{sporkSubRhsColor}{HTML}{006000}
\definecolor{sporkSubBoxColor}{HTML}{D8D8D8}
\newcommand{\sporkRule}[1]{
  %% {\color{sporkSubLhsColor} \setlength{\fboxrule}{0.9pt}\setlength{\fboxsep}{0.1pt} \boxed{\normalcolor \bullet} }
  %% \hspace{-2pt}
  %% \to
  %% \hspace{-2pt}
  %% {\color{sporkSubRhsColor} \setlength{\fboxrule}{0.9pt}\setlength{\fboxsep}{0.1pt} \boxed{\normalcolor \bullet} }
  \hyperref[sec:Semantics]{\hookrightarrow}
  \langle #1 \rangle
}
%% \sporkSub usage policy (internal; the reader only gets the definition in Semantics.tex):
%% * Box only what changes. With exceptions (below), unchanged context (operators, statement keywords, siblings)
%%   stays outside the box, even prefix/postfix context like -\sporkSub{..}{..} or \sporkSub{..}{..}:e.
%% * Box \sporksigma and \sporktr^* only in rules that can change them;
%%   an unboxed \sporksigma / \sporktr^* tells the reader the rule cannot change it.
%% * Substitution is purely syntactic and never reaches inside subscripts/superscripts:
%%   \sporktr^*_1 is a metavariable name, not \sporktr^* with something substituted.
%% * Either side may be empty (inserted or deleted syntax); write \sporkSubNothing for the empty side
%%   so the box doesn't collapse and the gap reads as intentional, e.g. \sporkSub{\sporkSubNothing}{\sporks'\texttt{;}}.
%% * Prefer duplicating a small unchanged neighbor over an empty side that deletes a suffix/prefix:
%%   strip annotations as \sporkSub{\sporkwz \# \sporkwinAnn}{\sporkwz}, not \sporkwz\sporkSub{\#\sporkwinAnn}{\sporkSubNothing}.
%%   Reserve empty sides for inserting/deleting whole syntax units (C-ForIter).
%% * Don't use empty sides to "unwrap" a child statement to a higher AST level
%%   (C-IfTrue, C-CudaWarps, C-CudaDeviceFunction): box the whole construct being replaced.
%%   Un-sequencing counts as lifting too: C-SeqDone boxes all of pass;s, not \sporkSub{pass;}{\sporkSubNothing}s.
%%   Inserting a sibling before an otherwise-unchanged node is fine (C-ForIter).
%% * Axioms that replace a construct with something unrelated (C-IfFalse -> pass)
%%   box the whole construct.
%% Em-width blank for an empty \sporkSub side.
\newcommand{\sporkSubNothing}{\quad}  % Redefine to nothing for the "in case of haters" scenario.
\newcommand{\sporkSub}[2]{
  {\color{sporkSubBoxColor} \setlength{\fboxrule}{0.75pt}\setlength{\fboxsep}{0.5pt} \boxed{
  {\color{sporkSubLhsColor} #1}
  \hspace{-2pt}
  \hyperref[sec:Semantics]{\color{lightgray}\blacktriangleright}
  \hspace{-2pt}
  {\color{sporkSubRhsColor} #2}}}
}

%% Convention: sporkSymbolName or sporksymbolName
%% Capitalization is based on camel case,
%% except the first character case must match the actual text case exactly.

%% Basic Grammar
\newcommand{\sporkbool}{\mathord{\hyperref[fig:BasicGrammar]{\textrm{bool}}}}
\newcommand{\sporkBool}{\mathord{\hyperref[fig:BasicGrammar]{\textrm{Bool}}}}
\newcommand{\sporkTrue}{\mathord{\hyperref[fig:BasicGrammar]{\texttt{True}}}}
\newcommand{\sporkFalse}{\mathord{\hyperref[fig:BasicGrammar]{\texttt{False}}}}
\newcommand{\sporkf}{\mathord{\hyperref[fig:BasicGrammar]{f}}}
\newcommand{\sporkF}{\mathord{\hyperref[fig:BasicGrammar]{\mathbf{F}}}}
\newcommand{\sporkfBehavior}{\mathord{\hyperref[fig:BasicGrammar]{f.\textsf{behavior}}}}
\newcommand{\sporkr}{\mathord{\hyperref[fig:BasicGrammar]{r}}}
\newcommand{\sporkv}{\mathord{\hyperref[fig:BasicGrammar]{v}}}
\newcommand{\sporkV}{\mathord{\hyperref[fig:BasicGrammar]{\mathbb{V}}}}
\newcommand{\sporkR}{\mathord{\hyperref[fig:BasicGrammar]{\mathbb{R}}}}
\newcommand{\sporkx}{\mathord{\hyperref[fig:BasicGrammar]{x}}}
\newcommand{\sporkX}{\mathord{\hyperref[fig:BasicGrammar]{\mathbb{X}}}}
\newcommand{\sporkextern}{\mathord{\hyperref[fig:BasicGrammar]{\textrm{extern}}}}
\newcommand{\sporkz}{\mathord{\hyperref[fig:BasicGrammar]{z}}}
\newcommand{\sporkZ}{\mathord{\hyperref[fig:BasicGrammar]{\mathbb{Z}}}}
\newcommand{\sporkmem}{\mathord{\hyperref[fig:BasicGrammar]{\textrm{mem}}}}
\newcommand{\sporkcontrol}{\mathord{\hyperref[fig:BasicGrammar]{\textrm{control}}}}
\newcommand{\sporkprim}{\mathord{\hyperref[fig:BasicGrammar]{\textrm{prim}}}}
\newcommand{\sporkspecialWindow}{\mathord{\hyperref[fig:BasicGrammar]{\textgpu{specialWindow}}}}
\newcommand{\sporkbarMech}{\mathord{\hyperref[fig:BasicGrammar]{\textgpu{barMech}}}}
\newcommand{\sporkmut}{\mathord{\hyperref[fig:BasicGrammar]{\textgpu{mut?}}}}
\newcommand{\sporkMut}{\mathord{\hyperref[fig:BasicGrammar]{\textgpu{Mut?}}}}
\newcommand{\sporkread}{\mathord{\hyperref[fig:BasicGrammar]{\textgpu{read}}}}
\newcommand{\sporkmutate}{\mathord{\hyperref[fig:BasicGrammar]{\textgpu{mutate}}}}
\newcommand{\sporksigma}{\mathord{\hyperref[fig:BasicGrammar]{\sigma}}}
\newcommand{\sporkSigma}{\mathord{\hyperref[fig:BasicGrammar]{\Sigma}}}
\newcommand{\sporktid}{\mathord{\hyperref[fig:BasicGrammar]{\textgpu{tid}}}}
\newcommand{\sporkTid}{\mathord{\hyperref[fig:BasicGrammar]{\textgpu{Tid}}}}
\newcommand{\sporktidMap}{\mathord{\hyperref[fig:BasicGrammar]{\textgpu{tidMap}}}}
\newcommand{\sporkTidMap}{\mathord{\hyperref[fig:BasicGrammar]{\textgpu{TidMap}}}}
\newcommand{\sporkloopMode}{\mathord{\hyperref[fig:BasicGrammar]{\textgpu{loopMode}}}}
\newcommand{\sporkLoopMode}{\mathord{\hyperref[fig:BasicGrammar]{\textgpu{LoopMode}}}}
\newcommand{\sporkcollUnit}{\mathord{\hyperref[fig:BasicGrammar]{\textgpu{collUnit}}}}
\newcommand{\sporkCollUnit}{\mathord{\hyperref[fig:BasicGrammar]{\textgpu{CollUnit}}}}
\newcommand{\sporkwarpName}{\mathord{\hyperref[fig:BasicGrammar]{\textgpu{warpName}}}}
\newcommand{\sporkwarpConfig}{\mathord{\hyperref[fig:BasicGrammar]{\textgpu{warpConfig}}}}

%% Window Grammar
\newcommand{\sporkqtl}{\mathord{\hyperref[fig:WindowGrammar]{\textgpu{qtl}}}}
\newcommand{\sporkQtl}{\mathord{\hyperref[fig:WindowGrammar]{\textgpu{Qtl}}}}
\newcommand{\sporkQtlPre}{\mathord{\hyperref[fig:WindowGrammar]{\underset{\textrm{pre}}{\textgpu{Qtl}}}}}
\newcommand{\sporkqtlInit}{\mathord{\hyperref[fig:WindowGrammar]{\underset{\textrm{init}}{\textgpu{qtl}}}}}
\newcommand{\sporkQtlAtom}{\mathord{\hyperref[fig:WindowGrammar]{\underset{\textrm{atom}}{\textgpu{Qtl}}}}}
\newcommand{\sporkQtlMask}{\mathord{\hyperref[fig:WindowGrammar]{\underset{\textrm{mask}}{\textgpu{Qtl}}}}}
\newcommand{\sporkQtlArrive}{\mathord{\hyperref[fig:WindowGrammar]{\underset{\textrm{Arrive}}{\textgpu{Qtl}}}}}
\newcommand{\sporkQtlAwait}{\mathord{\hyperref[fig:WindowGrammar]{\underset{\textrm{Await}}{\textgpu{Qtl}}}}}
\newcommand{\sporkc}{\mathord{\hyperref[fig:WindowGrammar]{{c}}}}
\newcommand{\sporkcz}{\mathord{\hyperref[fig:WindowGrammar]{\overset{z}{c}}}}
\newcommand{\sporkw}{\mathord{\hyperref[fig:WindowGrammar]{{w}}}}
\newcommand{\sporkwH}{\mathord{\hyperref[fig:WindowGrammar]{\overset{\#}{w}}}}
\newcommand{\sporkwz}{\mathord{\hyperref[fig:WindowGrammar]{\overset{z}{w}}}}
\newcommand{\sporkwzH}{\mathord{\hyperref[fig:WindowGrammar]{\overset{z\#}{w}}}}
\newcommand{\sporkwScalarH}{\mathord{\hyperref[fig:WindowGrammar]{\overset{\bullet\#}{w}}}}
\newcommand{\sporkwinAnn}{\mathord{\hyperref[fig:WindowGrammar]{\textgpu{winAnn}}}}
\newcommand{\sporkWinAnn}{\mathord{\hyperref[fig:WindowGrammar]{\textgpu{WinAnn}}}}

%% Expression Grammar
\newcommand{\sporkBinOp}{\mathord{\hyperref[fig:expr]{\textrm{BinOp}}}}
\newcommand{\sporke}{\mathord{\hyperref[fig:expr]{e}}}
\newcommand{\sporkeLo}{\mathord{\hyperref[fig:expr]{e_\textrm{lo}}}}
\newcommand{\sporkeHi}{\mathord{\hyperref[fig:expr]{e_\textrm{hi}}}}
\newcommand{\sporkExpr}{\mathord{\hyperref[fig:expr]{\textrm{Expr}}}}

%% Statement Grammar
\newcommand{\sporks}{\mathord{\hyperref[fig:stmt]{s}}}
\newcommand{\sporkStmt}{\mathord{\hyperref[fig:stmt]{\textrm{Stmt}}}}
\newcommand{\sporkMutateOpSet}{\hyperref[fig:stmt]{\{\texttt{=},\texttt{+=}\}}}

%% Procedure Grammar
\newcommand{\sporkformalParam}{\mathord{\hyperref[fig:ProcGrammar]{\textrm{formalParam}}}}
\newcommand{\sporkFormalParam}{\mathord{\hyperref[fig:ProcGrammar]{\textrm{FormalParam}}}}
\newcommand{\sporkproc}{\mathord{\hyperref[fig:ProcGrammar]{\textrm{proc}}}}
\newcommand{\sporkProc}{\mathord{\hyperref[fig:ProcGrammar]{\textrm{Proc}}}}

%% Unique Variable Names
\newcommand{\sporkalphaRename}{\mathord{\hyperref[sec:UniqueVariableNames]{\alpha\textrm{-rename}}}}
\newcommand{\sporkbinder}{\hyperref[sec:UniqueVariableNames]{\textbf{\textrm{($\sporkx$-binder)}}}}

%% Window Semantics
\newcommand{\sporkWindowSub}{\hyperref[fig:WindowSubSemantics]{\overset{\textrm{win}}{\mapsto}}}
\newcommand{\sporkWunroll}{\hyperref[fig:WindowUnrollSemantics]{\textrm{Wunroll}}}
\newcommand{\sporkWvalues}{\hyperref[fig:WindowUnrollSemantics]{\textrm{Wvalues}}}

%% Trace Grammar
\newcommand{\sporktr}{\mathord{\hyperref[fig:Traces]{\textgpu{tr}}}}
\newcommand{\sporkTr}{\mathord{\hyperref[fig:Traces]{\textgpu{Tr}}}}
\newcommand{\sporkTrNewAccess}{\mathord{\hyperref[fig:Traces]{\textgpu{TrNewAccess}}}}
\newcommand{\sporkTrFence}{\mathord{\hyperref[fig:Traces]{\textgpu{TrFence}}}}
\newcommand{\sporkTrArrive}{\mathord{\hyperref[fig:Traces]{\textgpu{TrArrive}}}}
\newcommand{\sporkTrAwait}{\mathord{\hyperref[fig:Traces]{\textgpu{TrAwait}}}}
\newcommand{\sporkTrReqCanAccess}{\mathord{\hyperref[fig:Traces]{\textgpu{TrReqCanAccess}}}}
\newcommand{\sporkTrReqPendingAwait}{\mathord{\hyperref[fig:Traces]{\textgpu{TrReqPendingAwait}}}}

%% Expr Semantics
\newcommand{\sporkERead}{\hyperref[fig:ExprSemantics]{\textsc{E-Read}}}

%% Data Semantics
\newcommand{\sporkDAssign}{\hyperref[fig:DataSemantics]{\textsc{D-Assign}}}
\newcommand{\sporkDReduce}{\hyperref[fig:DataSemantics]{\textsc{D-Reduce}}}
\newcommand{\sporkDInstrStepParam}{\hyperref[fig:DataSemantics]{\textsc{D-InstrStepParam}}}
\newcommand{\sporkDInstrStepBarrier}{\hyperref[fig:DataSemantics]{\textsc{D-InstrStepBarrier}}}
\newcommand{\sporkDInstrSub}{\hyperref[fig:DataSemantics]{\textsc{D-InstrSub}}}

%% Pending Await Grammar
\newcommand{\sporkpend}{\mathord{\hyperref[fig:PendingAwait]{\textgpu{pend}}}}
\newcommand{\sporkPend}{\mathord{\hyperref[fig:PendingAwait]{\textgpu{Pend}}}}
\newcommand{\sporkPendingAwait}{\mathord{\hyperref[fig:PendingAwait]{\textgpu{PendingAwait}}}}

%% SyncTraceHelpers
\newcommand{\sporktrAccessHelper}{\mathord{\hyperref[fig:SyncTraceHelpers]{\textgpu{trAccessHelper}}}}
\newcommand{\sporkinstrBarrierHelper}{\mathord{\hyperref[fig:SyncTraceHelpers]{\textgpu{instrBarrierHelper}}}}

%% SyncTraceSemantics
\newcommand{\sporkTrRead}{\hyperref[fig:SyncTraceSemantics]{\textsc{Tr-Read}}}
\newcommand{\sporkTrMutate}{\hyperref[fig:SyncTraceSemantics]{\textsc{Tr-Mutate}}}

%% Trace Checks
\newcommand{\sporkConcatRaw}{\mathord{\hyperref[sec:TraceChecks]{\underset{\textsf{raw}}{\texttt{++}}}}}
\newcommand{\sporkConcatChk}{\mathord{\hyperref[sec:TraceChecks]{\underset{\textsf{chk}}{\texttt{++}}}}}

%% Trace Vertices
\newcommand{\sporkvert}{\mathord{\hyperref[sec:TraceVertices]{\textsf{vert}}}}
\newcommand{\sporkVert}{\mathord{\hyperref[sec:TraceVertices]{\textsf{Vert}}}}
\newcommand{\sporktrVerts}{\mathord{\hyperref[eq:sporktrVerts]{\textsf{trVerts}}}}
\newcommand{\sporktrNewVerts}{\mathord{\hyperref[eq:sporktrNewVerts]{\textsf{trNewVerts}}}}
\newcommand{\sporktrNewVertsAt}{\mathord{\hyperref[eq:sporktrNewVertsAt]{\textsf{trNewVertsAt}}}}
\newcommand{\sporktlSig}{\mathord{\hyperref[sec:TraceVertices]{\textsf{tlSig}}}}
\newcommand{\sporkTlSig}{\mathord{\hyperref[sec:TraceVertices]{\textsf{TlSig}}}}
\newcommand{\sporkTlSigA}{\mathord{\hyperref[fig:VertexTable]{\textsf{TlSig1}}}}
\newcommand{\sporkTlSigB}{\mathord{\hyperref[fig:VertexTable]{\textsf{TlSig2}}}}
\newcommand{\sporkPendA}{\mathord{\hyperref[fig:VertexTable]{\textsf{Pend1}}}}
\newcommand{\sporkPendB}{\mathord{\hyperref[fig:VertexTable]{\textsf{Pend2}}}}
% VertexTable layout: one tikz diagram per trace event constructor.
% \vertexDiagram{constructor}{TlSig1}{TlSig2}{Pend1}{Pend2}
% The constructor box sits in the middle; TlSig boxes above, Pend boxes below,
% "1" boxes on the left, "2" boxes on the right. Constructor and values are math.
% An empty value means \varnothing, drawn as a minimized "label = \varnothing" box.
% Before drawing, each row (top/bottom) is measured:
%   * a box breaks after "=" only if its row would otherwise be wider than the
%     constructor box (wider box first; the other only if that makes the row fit;
%     \varnothing boxes never break);
%   * the row slides horizontally to stay within the constructor box's width
%     (centered on the constructor if it can't fit).
\makeatletter
\newsavebox\vd@tmpbox
\newdimen\vd@W \newdimen\vd@gap \vd@gap=3mm
\newdimen\vd@aOne \newdimen\vd@aTwo \newdimen\vd@bOne \newdimen\vd@bTwo
\newdimen\vd@a \newdimen\vd@b \newdimen\vd@lo \newdimen\vd@hi \newdimen\vd@tmp
\newdimen\vd@topShift \newdimen\vd@botShift
\newdimen\vd@cOne \newdimen\vd@cTwo \newdimen\vd@dOne \newdimen\vd@dTwo
\newif\ifvd@breakA \newif\ifvd@breakB
\newif\ifvd@breakTL \newif\ifvd@breakTR \newif\ifvd@breakBL \newif\ifvd@breakBR
\newcommand\vd@ifempty[3]{\if\relax\detokenize{#1}\relax #2\else #3\fi}
\newcommand\vd@fontfor[1]{\vd@ifempty{#1}{\scriptsize}{\small}}
\newcommand\vd@valfor[1]{\vd@ifempty{#1}{\varnothing}{#1}}
\tikzset{
  vertCtor/.style={draw, inner sep=3pt},
  vertVal/.style={draw, inner sep=2pt, align=left},
  vertEmpty/.style={draw, inner sep=1pt, align=left},
}
% \vd@widths{label}{value}{one-line width}{broken width}: outer node widths
\newcommand\vd@widths[4]{%
  \sbox\vd@tmpbox{\vd@fontfor{#2}$#1 = \vd@valfor{#2}$}#3=\wd\vd@tmpbox
  \sbox\vd@tmpbox{\vd@fontfor{#2}$#1 =$}#4=\wd\vd@tmpbox
  \sbox\vd@tmpbox{\vd@fontfor{#2}$\vd@valfor{#2}$}%
  \ifdim\wd\vd@tmpbox>#4 #4=\wd\vd@tmpbox\fi
  \vd@tmp=\vd@ifempty{#2}{2.4pt}{4.4pt}% 2 * inner sep + line width
  \advance#3 by \vd@tmp \advance#4 by \vd@tmp
  \ifdim#4>#3 #4=#3\fi
  \vd@ifempty{#2}{#4=#3}{}}% \varnothing boxes never break
\newcommand\vd@tryBreakA{\ifdim\vd@aTwo<\vd@aOne \vd@a=\vd@aTwo \vd@breakAtrue\fi}
\newcommand\vd@tryBreakB{\ifdim\vd@bTwo<\vd@bOne \vd@b=\vd@bTwo \vd@breakBtrue\fi}
\newif\ifvd@tooWide
\newcommand\vd@checkWidth{\vd@tmp=\vd@a \advance\vd@tmp by \vd@b \advance\vd@tmp by \vd@gap
  \ifdim\vd@tmp>\vd@W \vd@tooWidetrue \else \vd@tooWidefalse \fi}
% \vd@row{labelA}{valueA}{labelB}{valueB}{shift dimen}: sets \ifvd@breakA/B and the shift
\newcommand\vd@row[5]{%
  \vd@widths{#1}{#2}\vd@aOne\vd@aTwo \vd@widths{#3}{#4}\vd@bOne\vd@bTwo
  \vd@a=\vd@aOne \vd@b=\vd@bOne \vd@breakAfalse \vd@breakBfalse
  \vd@checkWidth
  \ifvd@tooWide
    \ifdim\vd@aOne<\vd@bOne \vd@tryBreakB \else \vd@tryBreakA \fi
    \vd@checkWidth
    \ifvd@tooWide % break the other box too, but only if that makes the row fit
      \vd@tmp=\vd@aTwo \advance\vd@tmp by \vd@bTwo \advance\vd@tmp by \vd@gap
      \ifdim\vd@tmp>\vd@W\else \vd@tryBreakA \vd@tryBreakB \fi
    \fi
  \fi
  \vd@lo=-0.5\vd@W \advance\vd@lo by 0.5\vd@gap \advance\vd@lo by \vd@a
  \vd@hi=0.5\vd@W \advance\vd@hi by -0.5\vd@gap \advance\vd@hi by -\vd@b
  \ifdim\vd@lo>\vd@hi #5=0.5\vd@lo \advance#5 by 0.5\vd@hi
  \else #5=0pt
    \ifdim#5<\vd@lo #5=\vd@lo\fi
    \ifdim#5>\vd@hi #5=\vd@hi\fi
  \fi}
% \vd@box{options}{node name}{label}{value}{break?}
\newcommand\vd@box[5]{%
  \vd@ifempty{#4}%
    {\node[vertEmpty, font=\scriptsize, #1] (#2) {$#3 = \varnothing$};}%
    {\node[vertVal, font=\small, #1] (#2) {#5{$#3 =$\\ $#4$}{$#3 = #4$}};}}
\newcommand\vd@ifTL{\ifvd@breakTL\expandafter\@firstoftwo\else\expandafter\@secondoftwo\fi}
\newcommand\vd@ifTR{\ifvd@breakTR\expandafter\@firstoftwo\else\expandafter\@secondoftwo\fi}
\newcommand\vd@ifBL{\ifvd@breakBL\expandafter\@firstoftwo\else\expandafter\@secondoftwo\fi}
\newcommand\vd@ifBR{\ifvd@breakBR\expandafter\@firstoftwo\else\expandafter\@secondoftwo\fi}
\newcommand{\vertexDiagram}[5]{%
  \sbox\vd@tmpbox{$\sporkvert = (\sporkz, #1)$}\vd@W=\wd\vd@tmpbox \advance\vd@W by 6.4pt
  \vd@row{\sporkTlSigA(\sporkvert)}{#2}{\sporkTlSigB(\sporkvert)}{#3}\vd@topShift
  \ifvd@breakA\vd@breakTLtrue\else\vd@breakTLfalse\fi
  \ifvd@breakB\vd@breakTRtrue\else\vd@breakTRfalse\fi
  \vd@row{\sporkPendA(\sporkvert)}{#4}{\sporkPendB(\sporkvert)}{#5}\vd@botShift
  \ifvd@breakA\vd@breakBLtrue\else\vd@breakBLfalse\fi
  \ifvd@breakB\vd@breakBRtrue\else\vd@breakBRfalse\fi
  \begin{tikzpicture}[baseline=(ctor.base)]
    \node[vertCtor] (ctor) {$\sporkvert = (\sporkz, #1)$};
    \vd@box{anchor=south east, at={($(ctor.north)+(\vd@topShift-0.5\vd@gap,2.5mm)$)}}{tl}{\sporkTlSigA(\sporkvert)}{#2}\vd@ifTL
    \vd@box{anchor=south west, at={($(ctor.north)+(\vd@topShift+0.5\vd@gap,2.5mm)$)}}{tr}{\sporkTlSigB(\sporkvert)}{#3}\vd@ifTR
    \vd@box{anchor=north east, at={($(ctor.south)+(\vd@botShift-0.5\vd@gap,-2.5mm)$)}}{bl}{\sporkPendA(\sporkvert)}{#4}\vd@ifBL
    \vd@box{anchor=north west, at={($(ctor.south)+(\vd@botShift+0.5\vd@gap,-2.5mm)$)}}{br}{\sporkPendB(\sporkvert)}{#5}\vd@ifBR
    \draw ($(tl.south east)+(-1.5mm,0)$) -- ++(0,-2.5mm);
    \draw ($(tr.south west)+(1.5mm,0)$) -- ++(0,-2.5mm);
    \draw ($(bl.north east)+(-1.5mm,0)$) -- ++(0,2.5mm);
    \draw ($(br.north west)+(1.5mm,0)$) -- ++(0,2.5mm);
  \end{tikzpicture}%
}
% \vertexDiagramSides: same arguments, but the "1" quadrants sit to the left of the
% constructor and the "2" quadrants to the right (TlSig upper, Pend lower).
% Trades horizontal space for vertical. If the diagram would exceed \linewidth,
% the wider column's boxes break after "=" (then the other column if still needed).
\newcommand\vd@sideWidths{%
  \vd@a=\ifvd@breakTL\vd@aTwo\else\vd@aOne\fi
  \vd@tmp=\ifvd@breakBL\vd@bTwo\else\vd@bOne\fi
  \ifdim\vd@tmp>\vd@a \vd@a=\vd@tmp\fi
  \vd@b=\ifvd@breakTR\vd@cTwo\else\vd@cOne\fi
  \vd@tmp=\ifvd@breakBR\vd@dTwo\else\vd@dOne\fi
  \ifdim\vd@tmp>\vd@b \vd@b=\vd@tmp\fi
  \vd@tmp=\vd@W \advance\vd@tmp by \vd@a \advance\vd@tmp by \vd@b
  \advance\vd@tmp by 5mm\relax}
\newcommand{\vertexDiagramSides}[5]{%
  \sbox\vd@tmpbox{$\sporkvert = (\sporkz, #1)$}\vd@W=\wd\vd@tmpbox \advance\vd@W by 6.4pt
  \vd@widths{\sporkTlSigA(\sporkvert)}{#2}\vd@aOne\vd@aTwo
  \vd@widths{\sporkPendA(\sporkvert)}{#4}\vd@bOne\vd@bTwo
  \vd@widths{\sporkTlSigB(\sporkvert)}{#3}\vd@cOne\vd@cTwo
  \vd@widths{\sporkPendB(\sporkvert)}{#5}\vd@dOne\vd@dTwo
  \vd@breakTLfalse \vd@breakBLfalse \vd@breakTRfalse \vd@breakBRfalse
  \vd@sideWidths
  \ifdim\vd@tmp>\linewidth
    \ifdim\vd@a<\vd@b \vd@breakTRtrue \vd@breakBRtrue \else \vd@breakTLtrue \vd@breakBLtrue \fi
    \vd@sideWidths
    \ifdim\vd@tmp>\linewidth \vd@breakTLtrue \vd@breakBLtrue \vd@breakTRtrue \vd@breakBRtrue \fi
  \fi
  \begin{tikzpicture}[baseline=(ctor.base)]
    \node[vertCtor] (ctor) {$\sporkvert = (\sporkz, #1)$};
    \vd@box{anchor=south east, at={($(ctor.west)+(-2.5mm,0.5mm)$)}}{tl}{\sporkTlSigA(\sporkvert)}{#2}\vd@ifTL
    \vd@box{anchor=south west, at={($(ctor.east)+(2.5mm,0.5mm)$)}}{tr}{\sporkTlSigB(\sporkvert)}{#3}\vd@ifTR
    \vd@box{anchor=north east, at={($(ctor.west)+(-2.5mm,-0.5mm)$)}}{bl}{\sporkPendA(\sporkvert)}{#4}\vd@ifBL
    \vd@box{anchor=north west, at={($(ctor.east)+(2.5mm,-0.5mm)$)}}{br}{\sporkPendB(\sporkvert)}{#5}\vd@ifBR
    \draw ($(tl.south east)+(0,1.2mm)$) -- ++(2.5mm,0);
    \draw ($(tr.south west)+(0,1.2mm)$) -- ++(-2.5mm,0);
    \draw ($(bl.north east)+(0,-1.2mm)$) -- ++(2.5mm,0);
    \draw ($(br.north west)+(0,-1.2mm)$) -- ++(-2.5mm,0);
  \end{tikzpicture}%
}
\makeatother
\newcommand{\sporkvertQtlMask}{\mathord{\hyperref[fig:VertexTable]{\textsf{vertQtlMask}}}}

%% TODO links to forthcoming definitions
\newcommand{\sporkConcatEnv}{\mathord{\underset{\textsf{env}}{\texttt{++}}}}
\newcommand{\sporkqtlConvergent}{\mathord{\textsf{qtlConvergent}}}
\newcommand{\sporkQtlCuda}{\mathord{\underset{\textrm{cuda}}{\sporkQtl}}}
